\documentclass[11pt]{article}
\usepackage[a4paper,margin=27mm]{geometry}
\usepackage{amsmath,amssymb,mathtools}
\usepackage{fontspec,unicode-math}
\setmainfont{Latin Modern Roman}
\setmathfont{Latin Modern Math}
\usepackage{xurl,hyperref}
\hypersetup{colorlinks=true,urlcolor=blue,linkcolor=blue}
\newcommand{\sref}[2]{\hyperlink{source:#1:#2}{[S#2]}}
\newcommand{\cataloguescope}[1]{\paragraph{Mathematical question.} #1}
\begin{document}
% UnsolvedMath ID 8700010; source catalogue code AMR-086-0010
\setcounter{section}{8700009}
\section{Conjecture 2.2 — Erdős-Woods}\label{8700010}
{\small\sffamily UnsolvedMath ID 8700010 · Graph Theory}
\subsection{Definitions and mathematical statement}

\paragraph{Shared notation.}$\mathbb N=\{1,2,\ldots\}$, and $\mathbb Z,\mathbb Q,\mathbb R,\mathbb C$ denote the integers, rationals, reals and complex numbers. Any other conventions are specified locally.


For a positive integer $a$, let $R(a)=\prod_{p\mid a}p$ be its radical, with $R(1)=1$.
\paragraph{Statement recorded in the supplied source.}
Does there exist an integer $k\ge1$ such that for every $m,n\ge1$, equality $R(m+i)=R(n+i)$ for all $0\le i<k$ implies $m=n$? \sref{8700010}{2}
\subsection{Short English statement}
Can a fixed number of successive prime-divisor sets uniquely determine the starting positive integer?
\subsection{Sources}
\begin{description}
\item[\hypertarget{source:8700010:1}{S1}] ULAM.AI and the UnsolvedMath Contributors, UnsolvedMath: A Curated Collection of Open Mathematics Problems, record 8700010 (AMR-086-0010). CC BY 4.0. \url{https://huggingface.co/datasets/ulamai/UnsolvedMath}.
\item[\hypertarget{source:8700010:2}{S2}] Michel Waldschmidt, Open Diophantine Problems, arXiv:math/0312440v2 (24 January 2004), Conjecture 2.2. \url{https://arxiv.org/abs/math/0312440}.
\end{description}
\label{end:8700010}
\end{document}
